\chapter{Architecture globale de la solution}

\section{Vue d'ensemble}

La solution est organisee autour de trois blocs independants mais connectes:
\begin{itemize}[leftmargin=*]
  \item \textbf{Backend} \texttt{football-academy}: API metier, securite JWT, temps reel WebSocket, stockage fichiers.
  \item \textbf{Frontend} \texttt{moez\_project}: application Flutter mobile + portail web admin.
  \item \textbf{Chatbot} \texttt{chatbot}: service Django/ML expose via endpoint HTTP local.
\end{itemize}

\section{Schema logique}

\begin{lstlisting}[style=codeblock]
+---------------------+            REST + STOMP WS            +------------------------+
|  Flutter Mobile App |  <----------------------------------> | Spring Boot Backend    |
|  (main.dart)        |                                        | /football-academy/api  |
+---------------------+                                        | /ws (STOMP)            |
          |
          | HTTP POST /api/chat
          v
+---------------------+
| Django Chatbot      |
| (localhost:8000)    |
| TF-IDF + fuzzy      |
+---------------------+

+---------------------+            REST + STOMP WS            +------------------------+
| Flutter Admin Web   |  <----------------------------------> | Spring Boot Backend    |
| (main_admin_web.dart)|                                       | (meme backend)         |
+---------------------+                                        +------------------------+
\end{lstlisting}

\section{Principe de separation des frontends}

Le frontend contient deux points d'entree distincts:
\begin{itemize}[leftmargin=*]
  \item \texttt{lib/main.dart} pour l'application mobile utilisateur;
  \item \texttt{lib/main\_admin\_web.dart} pour la plateforme admin web.
\end{itemize}

Cette approche permet de deployer les deux interfaces separement sans imposer les memes contraintes de distribution.

\section{Flux principaux}

\subsection*{Flux d'authentification}
\begin{enumerate}[leftmargin=*]
  \item l'utilisateur envoie ses credentials via l'app Flutter;
  \item le backend valide et emet un JWT;
  \item les routes sont protegees selon les roles (admin, trainer, etc.).
\end{enumerate}

\subsection*{Flux de chat temps reel}
\begin{enumerate}[leftmargin=*]
  \item le frontend charge conversations + contacts;
  \item le frontend s'abonne aux topics WebSocket;
  \item les nouveaux messages arrivent en push et mettent a jour UI.
\end{enumerate}

\subsection*{Flux chatbot}
\begin{enumerate}[leftmargin=*]
  \item l'utilisateur pose une question dans l'ecran chatbot Flutter;
  \item Flutter appelle \texttt{POST /api/chat} du service Django;
  \item Django retourne reponse + score + metadata de correspondance.
\end{enumerate}
